\scanpage{23}

\deleted{2\textsuperscript{o}}) \uline{Cas contraire}. Je suis pessimiste.

2\textsuperscript{o}) Cas $p=0$: Par exemple $k[[t]]$ sur $k$, ($k=\mathbb{Q}$, ou $k=\mathbb{C}\ldots$). Je présume que si $\dim A>0$, $A$ n'est jamais f\unsure{or\textsuperscript{t}|orm\textsuperscript{t}} simple sur $k$ pour sa topologie discrète.

Dans les cas défavorables, je doute qu'il y ait \deleted{\unsure{un|encore}} ``form\textsuperscript{t} simple \unsure{relativem\textsuperscript{t}|relativement|naturellement} au corps premier'' [Si $k$ n'est pas alg. sur le dit].

\[ \Omega^1_k\otimes_k A\longrightarrow\Omega^1_A\quad \text{admet un inverse à gauche?} \]

\deleted{Supposons $k$ de \unsure{degré|degré de} $\Omega_k^{\unsure{-|1|\vphantom{1}}}$} $k=\deleted{\unsure{k_0|k}}\mathbb{Q}((x_i)_{i\in I}\deleted{\unsure{i\in|i}})$ \deleted{$k=\mathbb{Q}$}

$A=k[[t]]$, $\Omega^1_k\otimes_k A\simeq A^{(I)}=\deleted{\unsure{\sum_{i\in I}|\bigoplus_{i\in I}}}\unsure{\sum|\bigoplus|\prod_i} A\, d_{/\mathbb{Q}}x_i$, on cherche $A\overset{D}{\unsure{-|\longrightarrow}}\Omega^1_k\otimes_k A$, qui soit donné sur $k$ par $D(x_i)=d_{/\mathbb{Q}}x_i$. \deleted{\unsure{Or|Car}} s'il y a une telle $D$, alors $D(t)$ est combinaison lin. finie de certains $d_{/\mathbb{Q}}x_i$, $i\in J$, et remplaçant $\mathbb{Q}$ par $k_0=\mathbb{Q}\cdot((x_i)_{i\in J})\deleted{\unsure{q|k}}$, on peut supposer \unsure{après|alors|que} $dt=0$. \deleted{$D\bigl(\sum a_it^i\bigr)=$}
% archive sheet number: 22
